% Resume: Arumalla Mohit Krishna
% Compile with pdfLaTeX.

\documentclass[letterpaper,10pt]{article}

\usepackage[margin=0.5in]{geometry}
\usepackage[T1]{fontenc}
\usepackage{lmodern}
\usepackage{microtype}
\usepackage{titlesec}
\usepackage{enumitem}
\usepackage[hidelinks]{hyperref}
\usepackage{tabularx}
\input{glyphtounicode}

\pdfgentounicode=1
\pagestyle{empty}
\urlstyle{same}
\raggedright
\setlength{\parindent}{0pt}
\setlength{\tabcolsep}{0pt}

\titleformat{\section}
  {\large\scshape}{}{0em}{}[\titlerule]

\titlespacing*{\section}{0pt}{6pt}{3pt}

\newcommand{\entry}[2]{%
  \noindent
  \begin{tabularx}{\linewidth}
    {@{}>{\raggedright\arraybackslash}X@{\hspace{6pt}}r@{}}
    \textbf{#1} & #2
  \end{tabularx}\par
}

\newcommand{\bullets}{
  \begin{itemize}[
    leftmargin=1.2em,
    topsep=2pt,
    itemsep=1pt,
    parsep=0pt,
    partopsep=0pt
  ]
}

\newcommand{\finishbullets}{
  \end{itemize}
}

\begin{document}

%----------HEADING----------
\begin{center}
    {\LARGE\bfseries Arumalla Mohit Krishna}\\[3pt]
    +1 (631) 352-8643 $|$
    \href{mailto:mohitkrishna.arumalla@stonybrook.edu}
    {mohitkrishna.arumalla@stonybrook.edu} $|$
    \href{https://linkedin.com/in/mohit-krishna-arumalla}
    {LinkedIn}
\end{center}
\vspace{-8pt}

%-----------EDUCATION-----------
\section{Education}

\entry
  {\textbf{Stony Brook University}}
  {\textbf{Aug 2025 -- Dec 2026}}

M.S. in Computer Science \hfill New York, USA\\
\textit{
Coursework: Natural Language Processing,
Introduction to Computer Vision, Vision Language Modeling
}
\vspace{3pt}

\entry
  {\textbf{Birla Institute of Technology and Science, Pilani}}
  {\textbf{Aug 2019 -- 2023}}

B.E. in Electronics \& Communication; GPA: 8.02/10 \hfill India
%-----------PROFESSIONAL SUMMARY-----------
\section{Professional Summary}

Software engineer with 2+ years of experience designing, building, and deploying \textbf{production-grade AI-powered software systems}. Currently pursuing an M.S. in Computer Science with hands-on expertise in \textbf{distributed services, semantic search, real-time data pipelines}, and \textbf{GPU model serving}. Skilled in rapid prototyping, integrating AI capabilities, and applying engineering best practices such as \textbf{code reviews, automated testing, CI/CD, observability, and security}. Proficient in \textbf{Python, Go, PyTorch}, and cloud infrastructure, with a strong focus on delivering measurable business impact through technology solutions.

%-----------TECHNICAL SKILLS-----------
\section{Technical Skills}

\textbf{Languages:}
Python, Go, C++, Java, SQL\\

\textbf{AI \& Machine Learning:}
PyTorch, TensorFlow/Keras, Transformers, scikit-learn,
Self-Supervised Learning, Vision Transformers,
Text Embeddings, Vector Search, mBART, BGE,
Parameter-Efficient Adaptation (LoRA), Visual Prompt Tuning\\

\textbf{Data \& Systems:}
gRPC, Kafka, Apache Flink, Flink SQL, Debezium/CDC,
PostgreSQL, Redis, Elasticsearch, Airflow, BigQuery\\

\textbf{Cloud \& Infrastructure:}
AWS, Docker, Kubernetes, NVIDIA Triton, Linux, Git,
CI/CD, Istio/Envoy, New Relic, Observability, Security Practices

%-----------EXPERIENCE-----------
\section{Experience}

\entry
  {\textbf{Software Engineer II $|$ Infra.Market, Bangalore}}
  {\textbf{Apr 2024 -- Aug 2025}}

\bullets

  \item Designed, built, and deployed a semantic product search system using Amazon Titan text embeddings and Elasticsearch vector search, enabling AI-driven natural-language product discovery through chat interfaces.

  \item Developed a high-throughput CDC pipeline using Flink, Kafka, Debezium, PostgreSQL, RocksDB, and Elasticsearch; reduced synchronization latency to ~500 ms with \textbf{exactly-once processing semantics} for reliable real-time data.

  \item Implemented Istio/Envoy service mesh across 15+ microservices with integrated New Relic distributed tracing for enhanced \textbf{monitoring, observability, and production debugging}.

\finishbullets

\entry
  {\textbf{Software Engineer I $|$ Infra.Market, Chennai}}
  {\textbf{Jul 2023 -- Mar 2024}}

\bullets

  \item Owned design and maintenance of Search and Location microservices, building gRPC services that processed 500K+ daily requests, optimizing AI-powered product discovery and geolocation capabilities.

  \item Resolved PostgreSQL connection leaks using Gorm ORM, reducing connection-pool exhaustion by 85\%; validated stability via concurrent load testing at 10K requests/second.

  \item Developed Airflow data pipelines for BigQuery to Elasticsearch migration, including query optimization and automated monitoring to ensure reliable and fresh search indices supporting customer-centric solutions.

\finishbullets

%-----------PROJECTS-----------
\section{Projects}

\entry
  {GPU-Optimized Personalized Product Search}
  {Aug 2026 -- Oct 2026}

\textit{Python, PyTorch, BGE, NVIDIA Triton, Go, gRPC}

\bullets

  \item Developed a two-tower retrieval model using \texttt{BAAI/bge-small-en-v1.5} query/product embeddings and trainable user tower for personalized AI search.

  \item Built and deployed an inference pipeline integrating precomputed product embeddings with approximate nearest-neighbor retrieval and a Go/gRPC gateway; incorporated NVIDIA Triton for efficient GPU query-encoder serving.

  \item Designed comprehensive benchmarks including latency (p95/p99), throughput, GPU utilization, and cost per search; evaluated retrieval metrics (Recall@K, NDCG@K) and system robustness under overload.

\finishbullets

\entry
  {Attention-Based Hallucination Detection in NMT}
  {Apr 2026 -- May 2026}

\textit{Python, PyTorch, mBART, Transformers, scikit-learn}

\bullets

  \item Created an English-German hallucination detection pipeline leveraging mBART attention features from WMT22/WMT23 quality estimation datasets.

  \item Extracted multi-head and multi-layer attention statistics including entropy, max mass, variance, and grounding ratios for effective representation of model confidence.

  \item Trained classifiers at token, sentence, and head levels achieving 0.712 F1 and 0.695 ROC-AUC; proposed hallucination-aware reranking methods integrating localization and risk signals.

\finishbullets

\entry
  {Prompt/LoRA for Self-Supervised Learning}
  {Ongoing}

\textit{Python, PyTorch, Vision Transformers}

\bullets

  \item Researching domain-specific visual prompting and LoRA within self-supervised ViT models to improve transferability and parameter efficiency across diverse image domains.

  \item Comparing inference performance with and without prompt/LoRA adaptations, designing evaluations for classification accuracy and robustness under domain shift.

\finishbullets

\end{document}